
\subsubsection{Key Findings}

\textbf{Methodology validation.} The checkpoint-gradient approach combined with capability detrending provides a framework for studying contamination-inflation relationships. The proof-of-concept validation demonstrates that the analysis pipeline can detect correlations when they exist in the data.

\textbf{Moderate effect size.} The observed correlation (r = 0.326) in simulated validation is weaker than the primary target (r > 0.5), suggesting that if contamination does inflate scores, capability may remain the primary driver of benchmark performance.

\textbf{Benchmark-specific effects.} Per-benchmark analysis shows that MMLU (r = 0.55) and ARC-Challenge (r = 0.53) exhibit stronger correlations than HellaSwag (r = 0.30) and WinoGrande (r = 0.25) in the simulated validation. This pattern is consistent with the hypothesis that factual benchmarks may be more susceptible to verbatim memorization effects.

\subsubsection{Limitations}

\textbf{Simulated validation only.} Current results derive from proof-of-concept validation using simulated contamination data where contamination-inflation correlation was structurally embedded in the data generation process. The correlation statistics demonstrate methodological feasibility; real Pythia checkpoint experiments are required to validate whether the observed relationship holds in practice.

\textbf{Resource requirements for full validation.} Full experimental validation requires:
\begin{itemize}
\item Approximately 144 GPU-hours for evaluation (72 checkpoints $\times$ ~2 hours each)
\item Approximately 48 CPU-hours for Pile n-gram index construction
\item Complete Pile corpus access (825 GB)
\end{itemize}

\textbf{Single model family.} The analysis focuses on Pythia models trained on The Pile. Generalization to other architectures and training corpora requires additional validation.

\textbf{Corpus coverage.} The 50,000 document subset represents 0.006\% of The Pile. Full corpus analysis would provide more comprehensive overlap statistics. Prior work \citep{yang2023rethinking} using full corpus analysis found 8-18\% overlap.

\textbf{Untested causal chain.} The methodology establishes correlation but does not verify the full causal mechanism (exposure leads to memorization leads to correct answers). The middle steps of the causal chain remain theoretical.

\textbf{Statistical significance variation.} Only 2 of 4 benchmarks (MMLU, ARC-Challenge) show statistically significant correlation at p < 0.05 in the per-benchmark analysis.

\subsubsection{Implications}

If contamination-inflation correlation is confirmed in real experiments, the methodology could enable:
\begin{itemize}
\item Contamination-aware benchmark score adjustment
\item Contamination-adjusted model comparison
\item Detection of adversarial training that maximizes benchmark scores through contamination
\end{itemize}

